\documentclass{scrartcl}
\usepackage{graphicx}  % required for inserting images

% for cyrillic and Bulgarian language
\usepackage[utf8]{inputenc}
\usepackage[T2A]{fontenc}
\usepackage[bulgarian]{babel}

% math packages
\usepackage{amsmath}
\usepackage{float}
\usepackage{amsfonts}
\usepackage{amssymb}
\usepackage{amsthm}
\usepackage{cancel}

\usepackage{ragged2e}  % adds FlushLeft
\usepackage{booktabs}
\usepackage{tabularx}
\usepackage{multirow}  % for complex tables

\usepackage{listings}
\usepackage{xcolor}

\definecolor{codegreen}{rgb}{0,0.6,0}
\definecolor{codegray}{rgb}{0.5,0.5,0.5}
\definecolor{codepurple}{rgb}{0.58,0,0.82}
\definecolor{backcolour}{rgb}{0.95,0.95,0.92}

\usepackage[margin=3cm]{geometry}

\lstdefinestyle{mystyle}{
    %caption=Example C++,
    label={lst:listing-cpp},
    language=C++,
    backgroundcolor=\color{backcolour},   
    commentstyle=\color{codegreen},
    keywordstyle=\color{magenta},
    numberstyle=\tiny\color{codegray},
    stringstyle=\color{codepurple},
    basicstyle=\ttfamily\footnotesize,
    breakatwhitespace=false,         
    breaklines=true,                 
    keepspaces=true,                 
    numbers=left,       
    numbersep=5pt,                  
    showspaces=false,                
    showstringspaces=false,
    showtabs=false,                  
    tabsize=2,
}
\lstset{style=mystyle}

\usepackage[colorlinks=true, linkcolor=blue, urlcolor=cyan]{hyperref}

\newtheorem{theorem}{Теорема}
\newtheorem{definition}{Дефиниция}

\title{Софтуерна архитектура. Проектиране и документиране на софтуерни архитектури}
\author{Ивайло Андреев + DeepSeek-V3}
\date{13 юни 2026}

\begin{document}
\maketitle
\tableofcontents
\newpage

\begin{FlushLeft}

\section{Понятие за софтуерна архитектура. Структури и изгледи. Необходимост. Влияние}

\subsection{Понятие за софтуерна архитектура}

Според Software Engineering Institute (SEI):

\begin{quote}
    \textit{“Софтуерната архитектура на дадена програма или изчислителна система е структурата или структурите на системата, които включват софтуерните елементи, външно видимите свойства на тези елементи и връзките между тях.”}
\end{quote}

С прости думи, софтуерната архитектура представлява:
\begin{itemize}
    \item \textbf{Фундаменталната организация} на системата.
    \item \textbf{Съвкупност от архитектурни структури} (структури и изгледи).
    \item \textbf{Елементи, връзки и ограничения} между тях.
    \item \textbf{Ключови проектни решения}, които оформят системата.
    \item \textbf{Основата за развитие} и поддръжка на системата през целия ѝ живот.
\end{itemize}

\subsection{Структури и изгледи (Structures and Views)}

Важно е да се прави разлика между \textbf{структура} и \textbf{изглед}:
\begin{itemize}
    \item \textbf{Структура (Structure):} Това е \textit{реалната организация} на системата – как елементите са свързани помежду си. Една система има много структури.
    \item \textbf{Изглед (View):} Това е \textit{представяне на дадена структура} за определена аудитория (например архитекти, разработчици, клиенти). Нито един изглед не описва цялата архитектура.
\end{itemize}

Архитектурните изгледи се делят на три основни групи:
\begin{enumerate}
    \item \textbf{Модулни изгледи (Module Views):}
    \begin{itemize}
        \item Декомпозиция на модулите.
        \item Употреба на модулите.
        \item Изглед на слоевете (многослоен).
        \item Йерархия на класовете.
    \end{itemize}
    \item \textbf{Изгледи на процесите (Process Views):} Показват изпълняваните процеси, нишки и тяхното взаимодействие.
    \item \textbf{Изгледи на разположението (Allocation/Deployment Views):}
    \begin{itemize}
        \item Изглед на внедряването (Deployment).
        \item Файлов изглед.
        \item Разпределение на работата.
    \end{itemize}
\end{enumerate}

\subsubsection{Модел 4+1}

Един популярен подход за представяне на архитектурата е моделът \textbf{4+1} (на Philippe Kruchten):
\begin{enumerate}
    \item \textbf{Логически изглед (Logical View):} Основните класове, компоненти и подсистеми.
    \item \textbf{Изглед на процесите (Process View):} Изпълняваните процеси, нишки и тяхното взаимодействие.
    \item \textbf{Изглед на кода (Development View):} Организацията на програмния код (модули, библиотеки, пакети).
    \item \textbf{Физически изглед (Deployment View):} Разполагането на софтуера върху хардуер (сървъри, клиенти, устройства).
    \item \textbf{(+1) Сценарии на употреба (Use Case/Scenarios):} Примери за реална работа на системата, които свързват останалите четири изгледа.
\end{enumerate}

\subsection{Необходимост от софтуерни архитектури}

Архитектурата е необходима поради няколко основни причини:
\begin{itemize}
    \item \textbf{Управление на сложността (Complexity Management):} Разделя системата на по-малки, разбираеми части. Позволява по-лесно разбиране на цялостната система.
    \item \textbf{Основа за анализ (Foundation for Analysis):} Служи като база за анализ на качеството, риска, цената и сроковете още в ранните етапи на разработката.
    \item \textbf{Комуникация (Communication):} Осигурява общ език и референтна рамка между различните заинтересовани страни – архитекти, разработчици, тестери, клиенти и мениджъри.
    \item \textbf{Повторна употреба (Reuse):} Позволява използването на доказани архитектурни шаблони и решения, което спестява време и ресурси.
\end{itemize}

\subsection{Влияние на архитектурата върху проекта и организацията (Architecture-Business Cycle)}

Връзката между бизнеса и архитектурата е двупосочна (цикъл):
\begin{itemize}
    \item \textbf{Бизнесът влияе върху архитектурата} чрез:
    \begin{itemize}
        \item \textbf{Изисквания} (функционални и качествени).
        \item \textbf{Бюджет} (ограничения върху ресурсите).
        \item \textbf{Срокове} (време за разработка).
    \end{itemize}
    \item \textbf{Архитектурата влияе върху бизнеса и организацията} чрез:
    \begin{itemize}
        \item \textbf{Структурата на екипа} (екипите често следват архитектурната модулна структура).
        \item \textbf{Управление на риска} (рисковете са идентифицирани на архитектурно ниво).
        \item \textbf{Разходите и поддръжката} (добрата архитектура намалява разходите за поддръжка и улеснява добавянето на нови функции).
    \end{itemize}
\end{itemize}

\newpage
\section{Архитектурни стилове}

\textbf{Архитектурен стил} е набор от:
\begin{itemize}
    \item \textbf{Компоненти} – изпълними единици, които извършват изчисления (напр. модули, класове, услуги).
    \item \textbf{Конектори} – механизми за комуникация между компонентите (напр. синхронни/асинхронни извиквания, съобщения).
    \item \textbf{Ограничения} – правила за това как компонентите могат да се комбинират.
    \item \textbf{Правила за взаимодействие} – семантика на комуникацията.
\end{itemize}

\subsection{Многослоен стил (Layered Architecture)}

При този стил системата се организира в йерархия от слоеве. \textbf{Всеки слой предоставя услуги само на слоя над себе си и използва услуги само от слоя под себе си.}

\paragraph{Предимства:}
\begin{itemize}
    \item \textbf{Разделяне на отговорностите} – всеки слой има ясна роля.
    \item \textbf{Скриване на информацията} – вътрешната реализация на слой е скрита.
    \item \textbf{Висока кохезия} – елементите в един слой са силно свързани функционално.
    \item \textbf{Лесна поддръжка и тестване} – слоевете могат да се тестват и заменят независимо.
\end{itemize}

\paragraph{Недостатъци:}
\begin{itemize}
    \item \textbf{Допълнителни зависимости} – понякога се налага „заобикаляне“ на слоевете, което нарушава модела.
    \item \textbf{По-ниска производителност} – заявката преминава през няколко слоя, което добавя забавяне.
\end{itemize}

\subsection{Модел-Изглед-Контролер (Model-View-Controller, MVC)}

MVC разделя системата на три взаимносвързани компонента:
\begin{enumerate}
    \item \textbf{Model (Модел):} Отговаря за данните и бизнес логиката. Уведомява View-тата при промяна.
    \item \textbf{View (Изглед):} Отговаря за визуализацията на данните (потребителския интерфейс).
    \item \textbf{Controller (Контролер):} Обработва входните действия от потребителя, променя модела и избира подходящ изглед.
\end{enumerate}

\paragraph{Предимства:}
\begin{itemize}
    \item \textbf{Разделяне на отговорностите} – ясно разделение на UI, логика и контрол.
    \item \textbf{Множество изгледи} – един и същ модел може да има различни изгледи.
    \item \textbf{Добра поддръжка} – промените в един компонент не влияят драстично на другите.
\end{itemize}

\paragraph{Недостатъци:}
\begin{itemize}
    \item \textbf{Допълнителна сложност} – за прости приложения MVC може да е излишен.
    \item \textbf{Допълнителен код} – изисква писане на повече код в сравнение с неструктуриран подход.
\end{itemize}

\subsection{Поточни архитектури (Pipe-and-Filter Architecture)}

Системата се състои от последователни компоненти, наречени \textbf{филтри (filters)}, които:
\begin{itemize}
    \item Получават входни данни.
    \item Обработват ги.
    \item Подават резултата към следващия компонент по \textbf{тръби (pipes)}.
\end{itemize}
\textit{Пример:} \texttt{Input → Filter1 → Filter2 → Filter3 → Output}

\paragraph{Предимства:}
\begin{itemize}
    \item \textbf{Повторна употреба} – филтрите могат да се пренареждат и използват в различни конфигурации.
    \item \textbf{Паралелизъм} – филтрите могат да работят паралелно, ако тръбите буферират данните.
    \item \textbf{Независими компоненти} – филтрите нямат знание един за друг.
\end{itemize}

\paragraph{Недостатъци:}
\begin{itemize}
    \item \textbf{Слаба производителност} – при много филтри и големи данни, серийната обработка може да е бавна.
    \item \textbf{Трудно споделяне на глобални данни} – филтрите са изолирани, което затруднява достъпа до споделено състояние.
\end{itemize}

\subsection{Архитектура, ориентирана към услуги (Service-Oriented Architecture, SOA)}

SOA е архитектурен стил, при който функционалността се предоставя като \textbf{услуги (services)} през мрежа. Основни участници:
\begin{itemize}
    \item \textbf{Доставчик на услуги (Service Provider):} Създава и публикува услуги.
    \item \textbf{Регистър на услуги (Service Registry):} Каталог, където услугите се регистрират и откриват.
    \item \textbf{Консуматор на услуги (Service Consumer):} Открива и използва услуги.
\end{itemize}

Основни операции в SOA:
\begin{enumerate}
    \item \textbf{Publish (Публикуване):} Доставчикът регистрира услугата в регистъра.
    \item \textbf{Find (Откриване):} Консуматорът търси услуга в регистъра.
    \item \textbf{Bind (Свързване):} Консуматорът се свързва към услугата и я използва.
\end{enumerate}

\paragraph{Характеристики:}
\begin{itemize}
    \item \textbf{Слаба свързаност (Loosely coupled):} Услугите са независими една от друга.
    \item \textbf{Повторна употреба (Reusability):} Услугите могат да се използват от множество приложения.
    \item \textbf{Оперативна съвместимост (Interoperability):} Услугите комуникират чрез стандартни протоколи (напр. HTTP, SOAP, REST).
\end{itemize}

\newpage
\section{Проектиране на софтуерната архитектура}

\subsection{Процес за проектиране}

Проектирането на архитектура е \textbf{цикличен процес}, който включва следните основни стъпки:
\begin{enumerate}
    \item \textbf{Анализ на изискванията:} Разбиране на функционалните и качествените изисквания, както и на бизнес ограниченията.
    \item \textbf{Избор на архитектурен стил:} Избор на подходящ стил (или комбинация от стилове) според анализа.
    \item \textbf{Декомпозиция:} Разделяне на системата на по-малки модули, компоненти или подсистеми.
    \item \textbf{Дефиниране на интерфейси:} Определяне на това как компонентите ще комуникират помежду си.
    \item \textbf{Оценка:} Проверка дали архитектурните решения удовлетворяват изискванията.
\end{enumerate}

\subsection{Архитектурни драйвери (Architectural Drivers)}

Това са основните фактори, които управляват архитектурните решения:
\begin{itemize}
    \item \textbf{Функционални изисквания:} Какво трябва да прави системата (напр. „системата трябва да генерира месечен отчет“).
    \item \textbf{Качествени показатели (Quality Attributes):} Производителност, сигурност, надеждност, използваемост, изменяемост и др.
    \item \textbf{Бизнес ограничения:} Срокове, бюджет, задължителни технологии, съответствие със стандарти.
\end{itemize}

\subsection{ADD (Attribute Driven Design)}

ADD е метод за проектиране на архитектурата, който поставя качествените показатели в центъра на процеса. Основните стъпки на ADD са:
\begin{enumerate}
    \item Избор на елемент за декомпозиция (например цялата система или голям модул).
    \item Избор на архитектурни драйвери за този елемент (кои качествени показатели са най-важни).
    \item Избор на архитектурен стил или шаблон, който адресира тези драйвери.
    \item Създаване на подмодули (декомпозиция на елемента).
    \item Дефиниране на интерфейсите на подмодулите.
    \item Проверка на решенията и повторение на процеса за подмодулите.
\end{enumerate}

\subsection{Тактики (архитектурни решения) за постигане на качествени показатели}

\textbf{Тактика} е архитектурно решение, което влияе директно върху постигането на даден качествен показател. Наборът от тактики образува \textbf{архитектурна стратегия}.

\subsubsection{Тактики за производителност (Performance)}
\begin{itemize}
    \item \textbf{Намаляване на изискванията:}
    \begin{itemize}
        \item По-добри алгоритми.
        \item Кеширане на данни.
        \item Намаляване на overhead (например използване на ефективни протоколи).
        \item Ограничаване на времето за изпълнение (таймаути).
    \end{itemize}
    \item \textbf{Управление на ресурсите:}
    \begin{itemize}
        \item Паралелизъм (многопоточност, разпределени изчисления).
        \item Излишък (redundancy) – кешове, реплики.
        \item Добавяне на повече ресурси (памет, процесорно време).
    \end{itemize}
    \item \textbf{Арбитраж (Scheduling):}
    \begin{itemize}
        \item FIFO (First-In-First-Out).
        \item Приоритети на заявките.
    \end{itemize}
\end{itemize}

\subsubsection{Тактики за улесняване на тестването (Testability)}
\begin{itemize}
    \item \textbf{Record \& Replay:} Записване на входни данни и повтаряне на тестовете.
    \item \textbf{Интерфейс срещу реализация (Interface vs Implementation):} Тестване чрез абстрактни интерфейси.
    \item \textbf{Специален тестов интерфейс (Test Interface):} Предоставяне на специални методи за тестване.
    \item \textbf{Мониторинг:} Наблюдение на състоянието на системата по време на тестове.
\end{itemize}

\subsubsection{Тактики за сигурност (Security)}
\begin{itemize}
    \item \textbf{Устояване на атаки:}
    \begin{itemize}
        \item Authentication (удостоверяване на самоличност).
        \item Authorization (оторизация – контрол на достъпа).
        \item Encryption (криптиране на данни).
        \item Integrity (проверка на целостта на данните).
    \end{itemize}
    \item \textbf{Възстановяване след атака:}
    \begin{itemize}
        \item Възстановяване на състоянието (restore).
        \item Одитен журнал (Audit trail) за проследяване на действията.
    \end{itemize}
\end{itemize}

\subsubsection{Тактики за удобство при употреба (Usability)}
\begin{itemize}
    \item \textbf{По време на изпълнение (за крайния потребител):}
    \begin{itemize}
        \item Обратна връзка (Feedback) за действията на потребителя.
        \item Отмяна (Undo) и анулиране (Cancel) на операции.
        \item Множество изгледи (Multiple Views) на данните.
        \item Баланс между автоматични действия и инициатива на потребителя.
    \end{itemize}
    \item \textbf{Насочени към разработчика на интерфейса:} API-та и инструменти за лесно изграждане на потребителски интерфейси.
\end{itemize}

\subsubsection{Тактики за изправност (Reliability)}
\begin{itemize}
    \item \textbf{Откриване и предпазване от откази:}
    \begin{itemize}
        \item Ping / Echo – проверка дали компонент е жив.
        \item Heartbeat – периодични сигнали за състояние.
        \item Exceptions – обработка на изключения.
    \end{itemize}
    \item \textbf{Отстраняване на откази:}
    \begin{itemize}
        \item Активен излишък (Active redundancy) – всички резервни компоненти работят паралелно.
        \item Пасивен излишък (Passive redundancy) – резервен компонент се активира при отказ.
        \item Извеждане от употреба на дефектния компонент.
        \item Следене на процесите (Process Monitoring).
    \end{itemize}
    \item \textbf{Повторно въвеждане в употреба:}
    \begin{itemize}
        \item Паралелна работа (shadow mode) – новата версия работи паралелно със старата.
        \item Контролни точки и връщане (Checkpoint / Rollback).
    \end{itemize}
\end{itemize}

\subsubsection{Тактики за изменяемост (Modifiability)}
\begin{itemize}
    \item \textbf{Локализиране на промените:}
    \begin{itemize}
        \item Поддържане на семантична свързаност (cohesion).
        \item Очакване на промени (предвиждане кои части вероятно ще се променят).
        \item Ограничаване на обхвата на опциите.
    \end{itemize}
    \item \textbf{Предотвратяване на ефекта на вълната (Ripple Effect):}
    \begin{itemize}
        \item Скриване на информация (Information Hiding).
        \item Ограничена комуникация между модулите.
        \item Използване на посредници (Mediators).
    \end{itemize}
    \item \textbf{Отлагане на свързването (Defer Binding):}
    \begin{itemize}
        \item Конфигурационни файлове.
        \item Компоненти с възможност за „plug-and-play“.
        \item Полиморфизъм и динамично зареждане.
    \end{itemize}
\end{itemize}

\newpage
\section{Документиране на софтуерната архитектура}

\subsection{Предназначение на документацията}

Архитектурната документация служи за:
\begin{itemize}
    \item \textbf{Комуникация} между всички заинтересовани страни.
    \item \textbf{Обучение} на нови членове на екипа.
    \item \textbf{Анализ} на архитектурата и вземане на информирани решения.
    \item \textbf{Поддръжка} и еволюция на системата.
    \item \textbf{Съхраняване на знания} за системата (knowledge management).
\end{itemize}

\subsection{Основен принцип на документиране}

\textbf{Архитектурата се документира чрез множество изгледи.} Нито един изглед не е достатъчен, за да опише цялата архитектура. Различните изгледи са насочени към различни аудитории и показват различни аспекти на системата.

\subsection{Съдържание и структура на документацията}

Една добра архитектурна документация трябва да съдържа следните елементи:

\paragraph{Основна информация:}
\begin{itemize}
    \item \textbf{Цели} на документа.
    \item \textbf{Контекст} на системата (как взаимодейства с външния свят).
    \item \textbf{Терминология} – дефиниции на основните понятия.
\end{itemize}

\paragraph{Архитектурни изгледи (Architectural Views):}
\begin{itemize}
    \item Модулни изгледи (например декомпозиция, слоеве).
    \item Изгледи на процесите (комуникация между процеси/нишки).
    \item Изгледи на разположението (хардуер, мрежа).
\end{itemize}

\paragraph{Отвъд изгледите (Beyond Views):}
\begin{itemize}
    \item \textbf{Архитектурни решения} – защо са взети определени решения.
    \item \textbf{Ограничения} – технологии, стандарти, бюджет.
    \item \textbf{Обосновка (Rationale)} – аргументите зад архитектурния дизайн.
\end{itemize}

\subsection{Документиране на архитектурен изглед}

Всеки архитектурен изглед трябва да съдържа следните елементи:
\begin{enumerate}
    \item \textbf{Основно представяне} – например UML диаграма.
    \item \textbf{Елементи и връзки} – описание на компонентите и конекторите.
    \item \textbf{Контекст} – как този изглед се свързва с останалите.
    \item \textbf{Варианти на използване} – как изгледът поддържа различни сценарии.
    \item \textbf{Архитектурна обосновка} – защо този изглед изглежда така.
    \item \textbf{Терминологичен речник} – специфични термини за изгледа.
    \item \textbf{Допълнителна информация} – препратки, версии.
\end{enumerate}

\subsection{Специфични архитектурни изгледи}

\subsubsection{Модулна декомпозиция (Module Decomposition View)}
\begin{itemize}
    \item Показва йерархичното разделяне на системата на модули.
    \item Връзките са от типа „X е под-модул на Y“.
    \item Всеки модул има ясно дефинирана отговорност, предоставя интерфейс и скрива вътрешната си реализация.
\end{itemize}

\subsubsection{Изглед на процесите (Process View)}
\begin{itemize}
    \item Елементите са процеси или нишки.
    \item Конекторите са комуникационни, синхронизационни или блокиращи операции.
    \item Показва паралелно изпълнение, времеви ограничения и комуникация между компонентите.
\end{itemize}

За представяне на изгледа на процесите най-често се използват следните UML диаграми:

\begin{table}[H]
\centering
\caption{Избор на UML диаграма за представяне на изглед на процесите}
\begin{tabular}{|p{3cm}|p{4cm}|p{5cm}|}
\hline
\textbf{Диаграма} & \textbf{Кога се използва} & \textbf{Какво показва} \\
\hline
Sequence Diagram & Когато има комуникация между процеси и ясен сценарий. & Последователност от съобщения между процеси във времето. \\
\hline
Activity Diagram & Когато има поток от действия и разклонения. & Последователност от действия и паралелни потоци. \\
\hline
State Machine Diagram & Когато системата има състояния и реагира на събития. & Състояния и преходи между тях. \\
\hline
\end{tabular}
\end{table}

\subsubsection{Изглед на разположението (Deployment View)}
\begin{itemize}
    \item Показва как софтуерните елементи се разполагат върху физическата среда (хардуер).
    \item Описва връзката между софтуер и хардуер.
    \item Представя комуникацията между устройства и среди.
\end{itemize}
Този изглед се създава след като вече имаме модулна декомпозиция и изглед на процесите, тъй като тогава компонентите са ясно дефинирани.

\section*{Заключение}

Софтуерната архитектура е фундаментът, върху който се изгражда всяка сложна софтуерна система. Тя включва структури и изгледи, които помагат за управление на сложността, комуникация и анализ. Различните архитектурни стилове – многослоен, MVC, поточен, SOA – предлагат доказани решения за различни проблемни области. Проектирането на архитектурата е цикличен процес, управляван от функционални и качествени изисквания, като тактиките за качествени показатели (производителност, сигурност, надеждност и др.) насочват конкретните архитектурни решения. Накрая, документацията чрез множество изгледи осигурява дълготрайна комуникация, поддръжка и еволюция на системата. Разбирането и правилното прилагане на тези концепции е от решаващо значение за успеха на всеки софтуерен проект.

\end{FlushLeft}
\end{document}